\documentclass[11pt]{article}
\usepackage[margin=1in]{geometry}
\usepackage{amsmath,amssymb,amsthm,mathtools}
\usepackage[T1]{fontenc}
\usepackage{lmodern,microtype,xurl}
\usepackage[colorlinks=true,linkcolor=blue,urlcolor=blue]{hyperref}
\newtheorem{theorem}{Theorem}[section]
\newtheorem{lemma}[theorem]{Lemma}
\newtheorem{proposition}[theorem]{Proposition}
\theoremstyle{remark}\newtheorem{remark}[theorem]{Remark}
\newcommand{\E}{\mathbb E}
\newcommand{\Prb}{\mathbb P}
\newcommand{\norm}[1]{\left\lVert #1\right\rVert}
\newcommand{\poly}{\operatorname{poly}}
\title{A positive-channel wire for deep softplus posteriors}
\author{ChatGPT-6 Astra \and Kaarel H\"anni \and Dmitry Vaintrob}
\date{26 October 2026}
\begin{document}
\maketitle
\begin{abstract}
For the standard softplus activation, a bounded-weight shallow Boolean
classifier can be transmitted through many hidden layers at isotropic
Gaussian energy $L+O(\log^2(LT))$, where $T$ is polynomial in the shallow
teacher resources. This matches the leading coefficient of the
nonconstant energy lower bound. The construction uses actual constant
neurons and charges their injection and final subtraction. Its operator
stability is polynomial even though the constant-channel weights depend
on a logarithmic anchor. Combined with a Lipschitz representation cover
and a Gaussian-distance posterior bound, it gives quasipolynomial
sample complexity at every finite learner width in cold limits and at
every fixed variance exponent $\gamma>1$, once depth exceeds a squared
logarithm of the teacher resources. The critical exponent $\gamma=1$,
arbitrary deep teachers, and the original tanh question remain outside
this result.
\end{abstract}

\section{Model, theorem, and proof dependencies}

Let $F(a)=\log(1+e^a)$, $X_d=\{-1,1\}^d$, and $d\ge1$. A learner has
$L\ge2$ matrices, $L-1$ softplus hidden layers of common width $N$,
augmented input $(1,x)$, no internal biases, and a linear output $z_W(x)$.
The independent Gaussian entry variances in the first matrix, internal
matrices, and readout are respectively
\[
 t^2/d,\qquad t^2/N,\qquad t^2/N^2,\qquad 0<t\le1.
\]
The scale may be prescribed as $t=t(d,L,N)$. For each fixed $\zeta>0$, the theorem is simultaneous
over $0<t\le N^{-\zeta/2}$, and therefore includes every such varying
choice, in particular $t(N)=N^{(1-\gamma)/2}$ with fixed
$\gamma=1+\zeta>1$.
For a target-labeled sample $S=((x_i,f(x_i)))_{i=1}^n$, set
\[
 E_S=\{|z_W(x)|>1/2\ \forall x\in X_d\}
 \cap\{\operatorname{sign}z_W(x_i)=f(x_i)\ \forall i\},
 \qquad Q_S=\mu(\,\cdot\mid E_S).
\]
The confidence-failure mass under uniform cube inputs is
\[
 e(S)=2^{-d}\#\{x:Q_S(\operatorname{sign}z_W(x)=f(x))\le2/3\}.
\]
Assume the target has a one-hidden-layer teacher of width $C$, physical
weights bounded by $B\ge1$, and signed output at least one everywhere.
Choose an integer
\[
 M_0=\left\lceil\max\{C,2,d+1,dC(d+1)B^2\}\right\rceil,
 \qquad T=M_0CB^2\ge1.
\]
These quantities are polynomial in the teacher resources.

\begin{theorem}\label{thm:main}
There are universal constants $c_0,c_1$ such that, if
\[
 L-1\ge c_0\log^2(T+2),
\]
then every width $N\ge C$ has the following learning guarantee. Put
$R=2+d+C+B+L+\varepsilon^{-1}+\log(2/\delta)$, with
$0<\varepsilon,\delta<1$.
Every limiting sign posterior as $t\downarrow0$, at a fixed finite
architecture, satisfies
\[
 n\ge\exp[c_1\log^2R]
 \quad\Longrightarrow\quad
 \Prb_S\{e(S)>\varepsilon\}\le\delta.
\]
For every fixed $\zeta>0$, the same holds for the actual finite-scale
posterior, with $c_1$ replaced by a constant $c_\zeta$, simultaneously for
all $N\ge C$ and $0<t\le N^{-\zeta/2}$. Thus the common isotropic
variance is at most $N^{-(1+\zeta)}$.
\end{theorem}

A limiting sign posterior means any subsequential limit on the finite
set of Boolean truth tables; the architecture is fixed before taking
that limit. The theorem is an upper bound for a shallow target class,
and does not bound arbitrary deep teachers by their width and depth.

We use one representation lemma from the logistic companion
\cite{logistic}: on a finite input set of radius at most two and
log-cardinality $O(d)$, globally margined networks with arbitrary
zero-fixing $1$-Lipschitz scalar activations, depth $L$, product of
Frobenius norms at most $P\ge1$, and final offset at most $2P$ have a
common finite sign class with log-cardinality $\poly(d,L,P)$.
Its proof compresses each layer by row-column importance sampling,
then stores finite scalar activation tables. It is independent of the
original width and requires no precision assumption on the original
weights or activations. Below we provide the softplus wire, its
stability, and the posterior transfer, including the small-width case.

\section{A matching-energy wire}

Use energy coordinates
\[
 A_1=\sqrt{d/N}\,W_1,\quad A_j=W_j\ (2\le j<L),\quad
 A_L=\sqrt N\,v,\qquad \kappa(A)=\sum_{j=1}^L\norm{A_j}_F^2.
\]
All coordinates have variance $\sigma^2=t^2/N$. The normalized input
is $(1,x)/\sqrt d$, hidden states are $h_j/\sqrt N$, and the normalized
activation is $F(\sqrt N a)/\sqrt N$. Since $F$ is $1$-Lipschitz, a
nonconstant globally absolute-margined classifier has, on some
opposite-label cube edge,
\begin{equation}\label{eq:lower}
 \prod_j\norm{A_j}_F>\sqrt d/2,\qquad
 \kappa(A)\ge L\exp[(\log d-\log4)/L]\ge L-\log4.
\end{equation}
This is a deterministic energy statement, not a finite-scale rate formula.

Suppose $N\ge4M_0$. Replicate every teacher row
$k=\lfloor N/(2M_0)\rfloor$ times, divide its outgoing coefficient by
$k$, and set all remaining first-layer rows and outgoing coefficients
to zero. The teacher computation $g=u\cdot h_1$ is unchanged. Its
first energy is at most two, and
\[
 D=N\norm u^2\le4M_0CB^2=4T.
\]
At least $N/2$ first-layer neurons remain with exact value
$q_0=F(0)=\log2$, on which $u$ vanishes. These supply a constant without
an internal bias.

Write $m=L-1\ge2$ and set
\[
 q=T^{1/(2m)},\qquad k_j=\min(j,m-j),\qquad
 r_j=\frac{q^j}{2\sqrt T\,k_j}\quad(1\le j\le m-1).
\]
Then $q^m=\sqrt T$, $r_j\le1/(2qk_j)\le1/2$, and
$r_{m-1}=1/(2q)$. Choose the positive anchor
\[
 A=4+3\log(m+2)+\log(T+2),\qquad
 a_j=\frac{r_j}{r_{j-1}}=q\frac{k_{j-1}}{k_j}\quad(2\le j\le m-1).
\]
The new hidden signal values are
\begin{equation}\label{eq:wire}
 h_1^{\rm new}=F(A+r_1g),\qquad
 h_j^{\rm new}=F\bigl(A+a_j(h_{j-1}^{\rm new}-A)\bigr).
\end{equation}
Implement each new layer using $n=\lfloor N/2\rfloor$ identical signal
neurons and $n$ zero-row constant neurons. Ignore an odd surplus neuron.
The coefficients into a signal neuron are $a_j$ on the preceding signal
average and $A(1-a_j)/q_0$ on the preceding constant average.
Thus every affine expression in \eqref{eq:wire} is implemented by actual
weights and the existing constant channel.

\begin{lemma}\label{lem:margin}
The final readout $z=4q(h_{m-1}^{\rm new}-A)$, implemented using the
signal and constant group averages, has signed margin at least one
whenever $|g|\ge1$.
\end{lemma}
\begin{proof}
At $g=\pm1$, compare the signal with $A+r_jg$. All additive errors are
nonnegative because $F(s)\ge s$, and all corresponding preactivations
are at least $A-1/2$. Each new additive error is therefore at most
$e^{-A+1/2}$. Consecutive ideal gains telescope:
\[
 \prod_{\ell=i+1}^j a_\ell=q^{j-i}\frac{k_i}{k_j}\le m\sqrt T.
\]
The final error is at most
\[
 m^2\sqrt T\,e^{-A+1/2}\le\frac1{8q}=\frac{r_{m-1}}4.
\]
The last inequality follows from $q\sqrt T\le T$ and the displayed
choice of $A$. Thus the final output is at least two at $g=1$ and
at most $-3/2$ at $g=-1$. Every map in the scalar chain and the final
readout is increasing in $g$. This extends the margin to all $|g|\ge1$;
an upper bound on the teacher output is not required here.
\end{proof}

\begin{proposition}\label{prop:energy}
If $m\ge c_0\log^2(T+2)$, the constructed witness has
\begin{equation}\label{eq:energy}
 \kappa(A_*)\le K:=L+C\log^2(LT+2)
\end{equation}
for a universal constant $C$.
\end{proposition}
\begin{proof}
The first new signal rows have disjoint teacher and constant supports.
Using at least $N/4$ reserved constants, their total energy is at most
\[
 (D/2)r_1^2+2A^2/q_0^2\le q^2/2+2A^2/q_0^2.
\]
Each later signal/constant layer has energy exactly
$a_j^2+(A^2/q_0^2)(1-a_j)^2$. The last readout has weighted energy
at most $64q^2(1+A^2/q_0^2)$, including the cost of subtracting $A$.
The factor $N/n\le3$ handles odd widths.

For $R_j=k_{j-1}/k_j$, the reciprocal pairs $k/(k+1)$ and $(k+1)/k$
give
\[
 \sum_jR_j^2\le m+C,\qquad \sum_j(1-R_j)^2\le C.
\]
Indeed the paired excess of squares above two is
$2/[k(k+1)]+k^{-2}+(k+1)^{-2}$, a summable sequence. Since $a_j=qR_j$,
\[
 \sum_j a_j^2\le(m+C)q^2,\qquad
 \sum_j(1-a_j)^2\le C\{m(q-1)^2+q^2\}.
\]
Including the first teacher layer yields
\[
 \kappa(A_*)\le mq^2+Cq^2+CA^2\{1+q^2+m(q-1)^2\}.
\]
The depth assumption bounds $q$ and $m(q-1)^2$ by universal constants,
while $mq^2\le m+C\log(T+2)$. Finally $A\le C\log(LT+2)$.
This proves \eqref{eq:energy}, with leading coefficient one as in
\eqref{eq:lower}. Every constant-channel contribution is charged.
\end{proof}

\section{Polynomial operator stability}

\begin{lemma}\label{lem:stability}
There is a polynomial $J=\poly(d,C,B,L)$, independent of $N\ge4M_0$,
such that perturbing every energy-coordinate matrix of the witness by
operator norm at most $J^{-1}$ preserves global signed margin $3/4$.
\end{lemma}

\begin{proof}
Let $s,c$ be the signal and constant group indicator vectors, each of
squared norm $n$. A new internal witness matrix is
\[
 W_j=s w_j^T,\qquad
 w_j=(a_j/n)s+[A(1-a_j)/(q_0n)]c.
\]
For every diagonal matrix $D_j$ with entries in $[0,1]$,
\[
 w_{j+1}^TD_js=\frac{a_{j+1}}n\sum_{i\in\mathrm{signal}}(D_j)_{ii}
 \in[0,a_{j+1}].
\]
Consequently a consecutive product of $D_jW_j$ has operator norm at most
\[
 \sqrt{a_i^2+[A(1-a_i)/q_0]^2}\prod_{\ell=i+1}^ja_\ell.
\]
The telescoping gains bound this by a polynomial in $L,T,A$.
The first injection, original teacher matrix, and final readout have
polynomial operator norms, so all relevant propagators are bounded by a
polynomial $P_0$. The bound permits arbitrary intervening diagonal
contractions, not merely derivatives at unperturbed states.

The shallow teacher range is polynomial in $d,C,B$. Also
\[
 |F(A+y)-A|\le |y|+\log2\qquad(A\ge0).
\]
This inequality and the telescoping gains give a polynomial bound $Q$
on every normalized witness state, uniformly on the cube and in $N$.
Write a perturbation block as $B_j$, with $\norm{B_j}_{\rm op}\le\eta$.
The componentwise mean-value formula, using $0<F'<1$, gives the exact
state-error identity
\[
 e_j=D_jA_j^*e_{j-1}+D_jB_j(u_{j-1}^*+e_{j-1}),
\]
where the $D_j$ are diagonal contractions. Variation of constants bounds
the maximum hidden-state error $E$ by
\[
 E\le\eta P_0L(Q+E).
\]
For $\eta P_0L\le1/2$, this gives $E\le2\eta P_0LQ$.
The final linear readout has polynomial norm, and its own perturbation
adds at most $\eta(Q+E)$. Taking a polynomially small $\eta$ makes total
output error at most $1/4$. Lemma~\ref{lem:margin} completes the proof.
\end{proof}

This uses the particular rank-one structure of the planted wire. No
polynomial partial-product bound is asserted for arbitrary networks of
energy $L+O(\log^2(LT))$.

\section{The actual posterior above the width threshold}

Fix an input anchor $x_0$. Center every normalized hidden state of a
realized network at that anchor. The resulting scalar activations are
\[
 a_{b,N}(r)=\frac{F(\sqrt N(b+r))-F(\sqrt N b)}{\sqrt N}.
\]
They fix zero and are $1$-Lipschitz. The linear matrices remain $A_j$,
and their product satisfies
\[
 P(A)\le(\kappa(A)/L)^{L/2}.
\]
The centered input has radius at most two. A nonconstant globally
margined output has an opposite-sign point, which bounds its anchor
offset by $2P(A)$. Constants can be listed separately. Therefore the
representation lemma from \cite{logistic} gives, at energy $K+1$,
a common finite sign class with
\begin{equation}\label{eq:class}
 P\le\exp[C\log^2(LT+2)],\qquad
 \log|\mathcal F|\le\exp[C\log^2(d+L+T+2)].
\end{equation}

For a finite-scale argument we must justify the posterior normalization,
not just compare witness energies. Set
$b^2=16\log9+8\log(8L)$, and let $\mathcal T$ be the intersection of
the parameter-block operator balls of radius $tb$. Since $N\ge4M_0$
and $M_0\ge d+1$, a unit-sphere net and Gaussian tails show
$\mu(\mathcal T)\ge1/2$. The set is closed, convex, and centrally
symmetric. Lemma~\ref{lem:stability} implies that, if $tb\le J^{-1}$,
then $A_*+\mathcal T\subset E_S$ for every target-labeled sample.
Gaussian translation and symmetry yield
\[
 \mu(E_S)\ge\tfrac12e^{-K/(2\sigma^2)}.
\]
Gaussian isoperimetry for a set of mass at least $1/2$ then gives
\begin{equation}\label{eq:distance}
 Q_S\{\operatorname{dist}(A,\mathcal T)^2>K+1\}
 \le2e^{-1/(2\sigma^2)}.
\end{equation}
These statements use the common isotropic variance $\sigma^2=t^2/N$.

On the complementary event write $A=U+B$ by projection onto $\mathcal T$,
so that $\kappa(U)\le K+1$ and each $B$ block has operator norm at most
$\eta=tb$. Center the \emph{realized} network $A$, freeze its scalar
maps $a_{b,N}$ and its actual anchor output, and replace only the linear
matrices by $U$. The surrogate is allowed to belong to the larger
Lipschitz representation class; its anchor maps are not recomputed, and
it need not be a native softplus network.

Zero-fixing Lipschitz telescoping and AM--GM bound the output change by
\begin{equation}\label{eq:perturb}
 2L\eta
 \left(\frac{(\sqrt{K+1}+\sqrt L\eta)^2}{L-1}\right)^{(L-1)/2}.
\end{equation}
Each of the $L$ terms is bounded by $2\eta$ times the other $L-1$
block norms, enlarged by at most $\eta$, and their squared sum is at
most $(\sqrt{K+1}+\sqrt L\eta)^2$. If
$\eta\le1/\sqrt{L(K+1)}$, this sum is at most $K+4$. Using
\eqref{eq:energy}, the power in \eqref{eq:perturb} is at most
$\exp[C\log^2(LT+2)]$. Thus a condition of the form
\begin{equation}\label{eq:small}
 tb\le\exp[-C\log^2(d+L+T+2)]
\end{equation}
both preserves the planted witness's margin and makes
\eqref{eq:perturb} at most $1/8$.

The surrogate has the same signs as the posterior network on the entire
cube, with margin greater than $3/8$. Its product obeys \eqref{eq:class}.
If it is nonconstant, its retained anchor offset is automatically bounded
by twice its own product, since its centered computation vanishes at
$x_0$. Hence the same sample-independent class $\mathcal F$ contains
posterior mass at least $1-2e^{-1/(2\sigma^2)}$.

For any fixed $\zeta>0$, the inequalities $t\le N^{-\zeta/2}$ and
$\sigma^2\le N^{-(1+\zeta)}$ ensure \eqref{eq:small} and
$2e^{-1/(2\sigma^2)}\le\beta$ above a threshold
\[
 N_*\le\exp\!\left[C_\zeta
   \log^2(2+d+C+B+L+\log(2/\beta))\right],\qquad N_*\ge4M_0.
\]
This is the large-width finite-scale part of the proof. It makes no
claim at the critical prefactor $t=1$.

\section{Small widths, cold limits, and learning}

For completeness, there is a lower-cost fallback for the purpose of
finite-width counting, although it does not match the coefficient in
\eqref{eq:lower}. The exact identity
\[
 F(g)-F(-g)=g
\]
provides a complementary relay. At any $N\ge\max(C,2)$, balance-clone
the first teacher rows to width $N$, dividing outgoing coefficients by
their replication counts. The first energy is at most
$2d(d+1)B^2$, and the weighted shallow readout energy is at most
$2C^2B^2$, independently of $N$. Split every added relay into two groups
of size $\lfloor N/2\rfloor$ and ignore an odd surplus neuron. Its
internal energy is four. The first relay costs at most the weighted
shallow readout energy, and the final readout costs at most six.
Thus a unit-margin depth-$L$ witness exists with energy polynomial in
$d,C,B,L$, independently of $N$. For $L=2$ use the cloned shallow witness
directly. If $N=1$, every represented truth table is constant because
the final hidden value is positive; one example determines its sign.

We spell out why the ordinary finite-width energy-net proof also applies
to this unbounded activation. Since $F(a)\le\log2+|a|$, on an energy
ball $\kappa\le R$ the normalized hidden states satisfy
\[
 \norm{u_j}\le\log2+\sqrt R\norm{u_{j-1}}.
\]
Consequently the logarithms of the state bounds and of a uniform
parameter-to-output Lipschitz constant $\Lambda_R$ are
$O(L\log(2+R+L))$, independent of $N$.
Let $K_0$ be the polynomial fallback witness energy, set $K'=2K_0+2$,
and take $r=\min(1,1/(4\Lambda_{K'}))$. A radius-$r$ witness neighborhood
retains signed margin $3/4$. In parameter dimension
$p=N(d+2)+(L-2)N^2$, use its inscribed cube with half-side
$\min(r/\sqrt p,\sigma)$. The cube mass gives
\[
 \mu(E_S)\ge e^{-K'/(2\sigma^2)-pT_0},\qquad
 T_0=\log\!\left(C\max\{1,\sqrt{p/N}/r\}\right).
\]
Clipping at $\sigma$ removes any divergence in $\log(1/t)$.
Dividing the elementary Gaussian energy tail by this denominator puts
posterior mass $1-\beta$ in the ball of radius squared
\[
 R_\beta=2K'+\frac4N
       \{p(T_0+\tfrac12\log2)+\log(1/\beta)\}.
\]
A net of mesh $1/(8\Lambda_{R_\beta+1})$ there preserves every globally
margined sign function. Its log-cardinality is polynomial in
$d,N,L,K_0,\log(1/\beta)$. Thus for $N<N_*$ it is quasipolynomial in
the resources of Theorem~\ref{thm:main}. Take the union of these finite
classes over the integers $C\le N<N_*$ and the common large-width class.
This adds at most $\log(N_*+1)$ to the largest log-cardinality bound,
which remains quasipolynomial. Thus one sample event covers every width
and eligible scale at the fixed depth; no width interval is omitted.

For a cold limit at fixed finite architecture, a strict-margin witness
of energy $K$ implies
\[
 Q_S\{\kappa>K+\rho\}\longrightarrow0\quad(\sigma\downarrow0)
 \qquad\text{for every }\rho>0.
\]
A fixed neighborhood retaining the witness's signs and energy at most
$K+\rho/2$ has Gaussian mass of order at least
$\sigma^{-p}e^{-(K+\rho/2)/(2\sigma^2)}$, while the tail above $K+\rho$
has the smaller exponential factor $e^{-(K+\rho)/(2\sigma^2)}$ times
a polynomial in $\sigma^{-1}$. For $N\ge4M_0$, the matching witness
and the class in \eqref{eq:class} therefore put every limiting sign
posterior on a common class of quasipolynomial log-cardinality.
The finite-width argument just given covers the remaining widths
uniformly in $t$ and hence in cold limits.

Finally, suppose a sample-independent finite class of size $M$ contains
posterior mass at least $1-\beta$. For
$n\ge(\log M+\log(1/\delta))/\alpha$, with probability at least
$1-\delta$ every sample-consistent member has true error at most $\alpha$.
The posterior Gibbs error is then at most $\alpha+\beta$. At a point
where correct-label confidence is at most $2/3$, wrong-label probability
is at least $1/3$, so $e(S)\le3(\alpha+\beta)$. Taking
$\alpha=\beta=\varepsilon/12$ proves the stated finite-scale bound.
For a limiting posterior supported on the class take $\beta=0$; passing
Gibbs error to the limit uses only continuity on the finite truth-table
simplex, not continuity of the inclusive confidence threshold.
This proves Theorem~\ref{thm:main}.

\section{A shallow parity witness and the scope limit}

Parity is an example within the shallow teacher class. Put
$k(x)=(d+\sum_i x_i)/2\in\{0,\ldots,d\}$. Full parity equals
$(-1)^d(-1)^k$, so a fixed readout sign supplies either convention.
The alternating values are
interpolated by
\[
 (-1)^k=1-2k+4\sum_{j=1}^{d-1}(-1)^{j+1}(k-j)_+
 \qquad(k\in\{0,\ldots,d\}).
\]
Replace each hinge by $F(\beta(k-j))/\beta$, whose uniform error is at
most $(\log2)/\beta$, and take $\beta=32d$. The accumulated error is
less than $1/8$. Realize the linear term exactly using
$F(k)-F(-k)=k$, and realize the constant using a zero-weight neuron of
value $\log2$. Multiplying the output by two gives signed margin at
least one with $d+2$ neurons and polynomially bounded weights.

When shallow teacher resources are polynomial in learner depth, the
depth condition in Theorem~\ref{thm:main} holds eventually, and its
quasipolynomial bound is subexponential in depth. Such shallow softplus
targets therefore do not supply exponential-depth lower bounds in the
cold or fixed-$\gamma>1$ regimes considered here. The proof does not
cover arbitrary deep teachers, $\gamma=1$, or arbitrary fixed positive
canonical prefactors. In particular, a both-label posterior cannot be
normalized by the mass of constant truth tables. These scope limits
also prevent any conclusion about the original exact-tanh problem.

\begin{thebibliography}{9}
\bibitem{logistic}
ChatGPT-6 Astra, K. H\"anni, and D. Vaintrob,
\emph{Shallow targets under deep logistic posteriors},
companion research proof, 2026. The finite Lipschitz representation
lemma and its layer-compression proof are used in Section 1.
\bibitem{borell}
C. Borell, \emph{The Brunn--Minkowski inequality in Gauss space},
Inventiones Mathematicae \textbf{30} (1975), 207--216.
\end{thebibliography}
\end{document}
